% Options for packages loaded elsewhere
\PassOptionsToPackage{unicode}{hyperref}
\PassOptionsToPackage{hyphens}{url}
\PassOptionsToPackage{dvipsnames,svgnames,x11names}{xcolor}
\documentclass[
  spanish,
  10pt,
  a4paper,
]{article}
\usepackage{xcolor}
\usepackage[top=2.2cm,bottom=2.2cm,left=2.1cm,right=2.1cm]{geometry}
\usepackage{amsmath,amssymb}
\setcounter{secnumdepth}{5}
\usepackage{iftex}
\ifPDFTeX
  \usepackage[T1]{fontenc}
  \usepackage[utf8]{inputenc}
  \usepackage{textcomp} % provide euro and other symbols
\else % if luatex or xetex
  \usepackage{unicode-math} % this also loads fontspec
  \defaultfontfeatures{Scale=MatchLowercase}
  \defaultfontfeatures[\rmfamily]{Ligatures=TeX,Scale=1}
\fi
\usepackage{lmodern}
\ifPDFTeX\else
  % xetex/luatex font selection
\fi
% Use upquote if available, for straight quotes in verbatim environments
\IfFileExists{upquote.sty}{\usepackage{upquote}}{}
\IfFileExists{microtype.sty}{% use microtype if available
  \usepackage[]{microtype}
  \UseMicrotypeSet[protrusion]{basicmath} % disable protrusion for tt fonts
}{}
\makeatletter
\@ifundefined{KOMAClassName}{% if non-KOMA class
  \IfFileExists{parskip.sty}{%
    \usepackage{parskip}
  }{% else
    \setlength{\parindent}{0pt}
    \setlength{\parskip}{6pt plus 2pt minus 1pt}}
}{% if KOMA class
  \KOMAoptions{parskip=half}}
\makeatother
\usepackage{color}
\usepackage{fancyvrb}
\newcommand{\VerbBar}{|}
\newcommand{\VERB}{\Verb[commandchars=\\\{\}]}
\DefineVerbatimEnvironment{Highlighting}{Verbatim}{commandchars=\\\{\}}
% Add ',fontsize=\small' for more characters per line
\newenvironment{Shaded}{}{}
\newcommand{\AlertTok}[1]{\textcolor[rgb]{1.00,0.00,0.00}{\textbf{#1}}}
\newcommand{\AnnotationTok}[1]{\textcolor[rgb]{0.38,0.63,0.69}{\textbf{\textit{#1}}}}
\newcommand{\AttributeTok}[1]{\textcolor[rgb]{0.49,0.56,0.16}{#1}}
\newcommand{\BaseNTok}[1]{\textcolor[rgb]{0.25,0.63,0.44}{#1}}
\newcommand{\BuiltInTok}[1]{\textcolor[rgb]{0.00,0.50,0.00}{#1}}
\newcommand{\CharTok}[1]{\textcolor[rgb]{0.25,0.44,0.63}{#1}}
\newcommand{\CommentTok}[1]{\textcolor[rgb]{0.38,0.63,0.69}{\textit{#1}}}
\newcommand{\CommentVarTok}[1]{\textcolor[rgb]{0.38,0.63,0.69}{\textbf{\textit{#1}}}}
\newcommand{\ConstantTok}[1]{\textcolor[rgb]{0.53,0.00,0.00}{#1}}
\newcommand{\ControlFlowTok}[1]{\textcolor[rgb]{0.00,0.44,0.13}{\textbf{#1}}}
\newcommand{\DataTypeTok}[1]{\textcolor[rgb]{0.56,0.13,0.00}{#1}}
\newcommand{\DecValTok}[1]{\textcolor[rgb]{0.25,0.63,0.44}{#1}}
\newcommand{\DocumentationTok}[1]{\textcolor[rgb]{0.73,0.13,0.13}{\textit{#1}}}
\newcommand{\ErrorTok}[1]{\textcolor[rgb]{1.00,0.00,0.00}{\textbf{#1}}}
\newcommand{\ExtensionTok}[1]{#1}
\newcommand{\FloatTok}[1]{\textcolor[rgb]{0.25,0.63,0.44}{#1}}
\newcommand{\FunctionTok}[1]{\textcolor[rgb]{0.02,0.16,0.49}{#1}}
\newcommand{\ImportTok}[1]{\textcolor[rgb]{0.00,0.50,0.00}{\textbf{#1}}}
\newcommand{\InformationTok}[1]{\textcolor[rgb]{0.38,0.63,0.69}{\textbf{\textit{#1}}}}
\newcommand{\KeywordTok}[1]{\textcolor[rgb]{0.00,0.44,0.13}{\textbf{#1}}}
\newcommand{\NormalTok}[1]{#1}
\newcommand{\OperatorTok}[1]{\textcolor[rgb]{0.40,0.40,0.40}{#1}}
\newcommand{\OtherTok}[1]{\textcolor[rgb]{0.00,0.44,0.13}{#1}}
\newcommand{\PreprocessorTok}[1]{\textcolor[rgb]{0.74,0.48,0.00}{#1}}
\newcommand{\RegionMarkerTok}[1]{#1}
\newcommand{\SpecialCharTok}[1]{\textcolor[rgb]{0.25,0.44,0.63}{#1}}
\newcommand{\SpecialStringTok}[1]{\textcolor[rgb]{0.73,0.40,0.53}{#1}}
\newcommand{\StringTok}[1]{\textcolor[rgb]{0.25,0.44,0.63}{#1}}
\newcommand{\VariableTok}[1]{\textcolor[rgb]{0.10,0.09,0.49}{#1}}
\newcommand{\VerbatimStringTok}[1]{\textcolor[rgb]{0.25,0.44,0.63}{#1}}
\newcommand{\WarningTok}[1]{\textcolor[rgb]{0.38,0.63,0.69}{\textbf{\textit{#1}}}}
\ifLuaTeX
\usepackage[bidi=basic,shorthands=off]{babel}
\else
\usepackage[bidi=default,shorthands=off]{babel}
\fi
\ifLuaTeX
  \usepackage{selnolig} % disable illegal ligatures
\fi
\setlength{\emergencystretch}{3em} % prevent overfull lines
\providecommand{\tightlist}{%
  \setlength{\itemsep}{0pt}\setlength{\parskip}{0pt}}
% Shared visual layer for the generated Kaleido FlowTwin reports.
\usepackage{helvet}
\renewcommand{\familydefault}{\sfdefault}
\usepackage{microtype}
\usepackage{xcolor}
\usepackage{titlesec}
\usepackage{fancyhdr}
\usepackage{enumitem}
\usepackage{booktabs}
\usepackage{xurl}
\usepackage{lastpage}
\usepackage{etoolbox}

\definecolor{KaleidoNavy}{HTML}{0A2342}
\definecolor{KaleidoBlue}{HTML}{176B87}
\definecolor{KaleidoCyan}{HTML}{3BB7C4}
\definecolor{KaleidoMint}{HTML}{DDF5F1}
\definecolor{KaleidoFog}{HTML}{F2F5F7}
\definecolor{KaleidoSlate}{HTML}{5D6A7A}
\definecolor{KaleidoCoral}{HTML}{F26B5B}

\titleformat{\section}
  {\Large\bfseries\color{KaleidoNavy}}
  {\thesection}{0.75em}{}[\vspace{0.2em}\color{KaleidoCyan}\titlerule]
\titleformat{\subsection}
  {\large\bfseries\color{KaleidoBlue}}
  {\thesubsection}{0.65em}{}
\titleformat{\subsubsection}
  {\normalsize\bfseries\color{KaleidoNavy}}
  {\thesubsubsection}{0.55em}{}
\titlespacing*{\section}{0pt}{2.2em}{1em}
\titlespacing*{\subsection}{0pt}{1.6em}{0.65em}

\pagestyle{fancy}
\fancyhf{}
\fancyhead[L]{\small\bfseries\color{KaleidoNavy}KALEIDO FLOWTWIN}
\fancyhead[R]{\small\color{KaleidoSlate}DOSSIER TECNICO}
\fancyfoot[L]{\scriptsize\color{KaleidoSlate}EVOCON Solutions GmbH}
\fancyfoot[R]{\scriptsize\color{KaleidoSlate}\thepage/\pageref*{LastPage}}
\renewcommand{\headrulewidth}{0.6pt}
\renewcommand{\headrule}{\hbox to\headwidth{\color{KaleidoCyan}\leaders\hrule height \headrulewidth\hfill}}
\renewcommand{\footrulewidth}{0pt}
\setlength{\headheight}{20pt}

\setlist[itemize]{leftmargin=1.4em,itemsep=0.25em,topsep=0.35em}
\setlist[enumerate]{leftmargin=1.7em,itemsep=0.25em,topsep=0.35em}
\setlength{\parindent}{0pt}
\setlength{\parskip}{0.55em}
\setlength{\emergencystretch}{3em}
\renewcommand{\arraystretch}{1.18}

\AtBeginEnvironment{quote}{%
  \color{KaleidoNavy}\itshape
  \setlength{\leftskip}{1em}
  \setlength{\rightskip}{1em}
}

\makeatletter
\renewcommand{\maketitle}{%
  \hypersetup{pageanchor=false}
  \begin{titlepage}
    \pagecolor{KaleidoNavy}
    \color{white}
    \vspace*{1.4cm}
    {\color{KaleidoCyan}\rule{2.4cm}{2.5pt}\par}
    \vspace{1.1cm}
    {\Huge\bfseries\@title\par}
    \vspace{0.65cm}
    {\Large\color{KaleidoMint}Kaleido FlowTwin\par}
    \vfill
    {\large Álvaro Schwiedop Souto\par}
    \vspace{0.15cm}
    {\normalsize\color{KaleidoMint}Industrial AI \& Predictive Quality Lead\par}
    \vspace{0.15cm}
    {\normalsize\color{KaleidoMint}EVOCON Solutions GmbH\par}
    \vspace{0.25cm}
    {\normalsize\color{KaleidoMint}\@date\par}
  \end{titlepage}
  \hypersetup{pageanchor=true}
  \nopagecolor
  \color{black}
}
\makeatother
\usepackage{bookmark}
\IfFileExists{xurl.sty}{\usepackage{xurl}}{} % add URL line breaks if available
\urlstyle{same}
% fallback for those not using the hyperref driver hyperxmp:
\makeatletter
\@ifundefined{xmpquote}{\newcommand{\xmpquote}[1]{#1}}{}
\makeatother
\hypersetup{
  pdftitle={Kaleido FlowTwin},
  pdflang={es},
  colorlinks=true,
  linkcolor={KaleidoBlue},
  filecolor={Maroon},
  citecolor={Blue},
  urlcolor={KaleidoBlue},
  pdfcreator={LaTeX via pandoc}}

\title{Kaleido FlowTwin}
\author{}
\date{15 julio 2026}

\begin{document}
\maketitle

{
\setcounter{tocdepth}{3}
\tableofcontents
}
\section{Kaleido FlowTwin}\label{kaleido-flowtwin}

Especificacion de un MVP de inteligencia operativa para Kaleido Tech. El
nombre \texttt{FlowTwin} es provisional y no presupone aprobacion de
marca.

\subsection{Veredicto ejecutivo}\label{veredicto-ejecutivo}

Ninguno de los repositorios existentes encaja de extremo a extremo con
el caso de Kaleido. El punto de partida mas solido es reutilizar la
gobernanza, los contratos de datos y los gates de
\texttt{predictiveops-worldmodel}; el modelo y los conectores deben ser
nuevos y orientados a eventos logisticos.

La oportunidad recomendada es una capa predictiva que amplie los
productos ya existentes de Kaleido:

\begin{itemize}
\tightlist
\item
  \texttt{Trace\ Port}: riesgo de desviacion, tiempo restante y cuello
  de botella por proyecto, operacion y turno.
\item
  \texttt{Shipping\ Board}: riesgo de retraso de expediciones y
  recepciones.
\item
  \texttt{Freight\ Intelligence}: riesgo de excepcion y ETA
  probabilistica.
\end{itemize}

El primer piloto debe ejecutarse sobre \textbf{una operacion repetible
de Trace Port}. No se promete entrenar un world model si el historico no
lo permite. El piloto produce valor en una escalera de evidencia:

\begin{enumerate}
\def\labelenumi{\arabic{enumi}.}
\tightlist
\item
  auditoria y contrato de datos;
\item
  process mining y conformance checking;
\item
  baselines de tiempo restante y riesgo;
\item
  simulacion discreta para escenarios;
\item
  modelo secuencial o JEPA solo si mejora los baselines en holdout.
\end{enumerate}

\subsection{Pregunta de negocio}\label{pregunta-de-negocio}

\begin{quote}
Dado el estado actual de una carga, descarga o manipulacion portuaria,
el plan, los recursos y el contexto, ¿cual es la probabilidad de
desviacion en las proximas 2/4/8 horas, cuanto falta para terminar,
donde esta el cuello de botella y que accion permitida reduce el riesgo?
\end{quote}

La salida es asesoramiento read-only. No escribe en TOS, ERP, PLC,
equipos ni sistemas de seguridad.

\subsection{Por que encaja con
Kaleido}\label{por-que-encaja-con-kaleido}

Kaleido ya digitaliza la operacion mediante Trace Port, Shipping Board y
Freight Intelligence. FlowTwin convierte los eventos que esas
herramientas generan en una capacidad adicional:

\begin{itemize}
\tightlist
\item
  pasar de visibilidad a anticipacion;
\item
  reducir supervision manual y deteccion tardia;
\item
  cuantificar lead time y falsas alarmas;
\item
  ofrecer un modulo premium a terminales y clientes;
\item
  mejorar estimacion, planificacion y revision post-operacion;
\item
  crear un activo de datos acumulativo sin exigir una integracion
  invasiva.
\end{itemize}

Fuentes corporativas auditadas a 15-07-2026:

\begin{itemize}
\tightlist
\item
  \url{https://www.kaleidologistics.com/es/kaleido-tech/}
\item
  \url{https://www.kaleidologistics.com/en/productos-ktech/trace-port/}
\item
  \url{https://www.kaleidologistics.com/en/productos-ktech/shipping-board/}
\item
  \url{https://www.kaleidologistics.com/en/productos-ktech/freight-intellligence/}
\item
  \url{https://www.kaleidologistics.com/en/port-operations/}
\end{itemize}

\subsection{Alcance del MVP}\label{alcance-del-mvp}

\subsubsection{Entrada minima}\label{entrada-minima}

\begin{itemize}
\tightlist
\item
  identificador de proyecto/operacion/turno;
\item
  plan previsto y milestones;
\item
  eventos con timestamp;
\item
  unidades de carga o lineas de packing list;
\item
  equipo/recurso/rol cuando exista;
\item
  inicio, pausa, incidencia y fin;
\item
  contexto: tipo de carga, buque, muelle, clima y turno;
\item
  resultado: duracion final, desviacion e incidencias;
\item
  costes solo si Kaleido quiere traducir riesgo a euros.
\end{itemize}

\subsubsection{Salidas}\label{salidas}

\begin{itemize}
\tightlist
\item
  tiempo restante con intervalo;
\item
  riesgo de exceder plan en 2/4/8 horas;
\item
  riesgo de incidencia, solo si hay etiquetas suficientes;
\item
  siguiente estado/evento probable;
\item
  cuello de botella y factores asociados;
\item
  nivel de confianza y abstencion;
\item
  escenarios what-if dentro de acciones aprobadas;
\item
  dashboard y export JSON/CSV/PDF.
\end{itemize}

\subsubsection{No incluido en la primera
version}\label{no-incluido-en-la-primera-version}

\begin{itemize}
\tightlist
\item
  control automatico;
\item
  vision generativa o video sintetico;
\item
  gemelo fisico 3D del puerto;
\item
  recomendacion causal sin acciones verificadas;
\item
  entrenamiento de un foundation model desde cero;
\item
  afirmaciones SOTA o ROI no medido.
\end{itemize}

\subsection{Documentos}\label{documentos}

\begin{itemize}
\tightlist
\item
  \href{docs/MEETING_BRIEF.md}{MEETING\_BRIEF.md}: guion completo para
  la reunion.
\item
  \href{docs/investigacion/KALEIDO_RESEARCH.md}{KALEIDO\_RESEARCH.md}:
  labor, productos y encaje de Kaleido.
\item
  \href{docs/investigacion/REPO_REUSE_AUDIT.md}{REPO\_REUSE\_AUDIT.md}:
  que reutilizar y que no.
\item
  \href{docs/investigacion/SOTA_2026.md}{SOTA\_2026.md}: sintesis
  tecnica y decisiones derivadas.
\item
  \href{docs/investigacion/REFERENCES.bib}{REFERENCES.bib}: bibliografia
  de trabajo.
\item
  \url{ARCHITECTURE.md}: arquitectura logica y contratos.
\item
  \url{DATASETS.md}: fuentes internas, publicas, licencias y uso
  permitido.
\item
  \url{DATA_REQUEST.md}: peticion minima para evaluar el encaje.
\item
  \url{PLAN.md}: hitos, gates y criterios de aceptacion.
\item
  \url{AGENTS.md}: contrato de implementacion para un agente.
\item
  \href{presentacion/Kaleido_FlowTwin_Presentacion.tex}{Kaleido\_FlowTwin\_Presentacion.tex}:
  fuente Beamer.
\item
  \href{presentacion/Kaleido_FlowTwin_Presentacion.pdf}{Kaleido\_FlowTwin\_Presentacion.pdf}:
  presentacion PDF.
\item
  \href{presentacion/Kaleido_FlowTwin_Presentacion.html}{Kaleido\_FlowTwin\_Presentacion.html}:
  presentacion web offline.
\item
  \href{correspondencia/CorreoEnviado2.txt}{CorreoEnviado2.txt}:
  respuesta propuesta a Manuel.
\end{itemize}

\subsection{Estructura}\label{estructura}

\begin{Shaded}
\begin{Highlighting}[]
\NormalTok{Kaleido Project/}
\NormalTok{|{-}{-} README.md, AGENTS.md, PLAN.md        \# entrada y contrato de desarrollo}
\NormalTok{|{-}{-} ARCHITECTURE.md, DATASETS.md         \# especificacion tecnica}
\NormalTok{|{-}{-} DATA\_REQUEST.md                      \# peticion minima a Kaleido}
\NormalTok{|{-}{-} correspondencia/                     \# correos y adjunto original}
\NormalTok{|{-}{-} presentacion/                        \# PDF, HTML y fuente TeX}
\NormalTok{\textasciigrave{}{-}{-} docs/                                \# guion e investigacion}
\end{Highlighting}
\end{Shaded}

\subsection{Estado}\label{estado}

\texttt{proposal\_ready}, sin datos internos de Kaleido y sin resultados
de modelo.

Todo numero de rendimiento mostrado en una futura reunion debe proceder
de un artefacto generado sobre un split congelado. Los resultados
historicos invalidados de otros repositorios no se pueden reutilizar
como prueba.

\end{document}
